\week{W4}{Parallel manipulators: mobility, loop closure, four-bar, delta}{4}
% W4 (booklet v3, condensed). Numbers verified by vcheck/v2_w04.py (and site/src/checks/w04.py).
\begin{multicols}{3}

\hd{Serial vs parallel; G--K reminder}
\begin{fbx}[Serial vs parallel \pg{74}]
\begin{tabularx}{\linewidth}{@{}p{0.27\linewidth}XX@{}}\toprule
 & \textbf{Serial} & \textbf{Parallel}\\\midrule
Chains & open & closed\\
Joints & all active & active + passive\\
FK; IK & simple; complex & \textbf{complex}; \textbf{simple}\\
Joint errors & accumulate & do not\\
Work volume & large & small\\
Payload/weight; inertia & low; high & very high; low\\\bottomrule
\end{tabularx}
\end{fbx}
G--K: $M=D(N-J-1)+\sum F_i$; $D=3$ planar, 6 spatial; $N$ counts the ground; R, P: 1; U: 2; S: 3. Full recipe and joint table: W1, p.\,\pageref{p:W1}.
\begin{trap}
``Parallel FK is easy'': wrong way round. Planar mechanism with $D=6$ (slider-crank gives $-2$, not 1). \chk
\end{trap}
\textbf{Workspace} \pg{21--23}: reachable (any orientation), dexterous (all orientations), constant-orientation. Parallel: each leg's $[\ell_{min},\ell_{max}]$ gives an annulus/shell about its base pivot; the workspace is their \textbf{intersection}. Planar 2 legs, pivots $(0,0),(0.6,0)$, $0.3\le\ell\le0.7$: apex $(0.30,0.632)$, area $y\ge0$ about $0.175\ \mathrm{m^2}$. \chk
\begin{center}
\begin{tikzpicture}[x=3.2cm,y=3.2cm,line width=0.5pt]
\begin{scope}
 \clip (-0.75,0) rectangle (1.35,0.8);
 \begin{scope}[even odd rule]\clip (0,0) circle (0.7) (0,0) circle (0.3);
 \begin{scope}[even odd rule]\clip (0.6,0) circle (0.7) (0.6,0) circle (0.3);
 \fill[g2] (-0.8,0) rectangle (1.4,0.8);
 \end{scope}\end{scope}
 \draw[dashed] (0,0) circle (0.7); \draw (0,0) circle (0.3);
 \draw[dashed] (0.6,0) circle (0.7); \draw (0.6,0) circle (0.3);
\end{scope}
\draw[line width=0.8pt] (-0.75,0)--(1.35,0);
\fill (0,0) circle (0.016) node[below=1pt,font=\scriptsize]{$\mathbf b_1$};
\fill (0.6,0) circle (0.016) node[below=1pt,font=\scriptsize]{$\mathbf b_2$};
\fill (0.3,0.6325) circle (0.016) node[above,font=\scriptsize]{$(0.30,0.632)$};
\fill (0.6333,0.2981) circle (0.016) node[right,font=\scriptsize]{$(0.633,0.298)$};
\fill (-0.0333,0.2981) circle (0.016) node[left,font=\scriptsize]{$(-0.033,0.298)$};
\node[font=\scriptsize] at (0.3,0.3) {workspace};
\end{tikzpicture}
\end{center}
\drawncap{solid: $\ell_{min}$; dashed: $\ell_{max}$; mirror below the axis if collisions ignored.}

\hd{Extended Gr\"ubler--Kutzbach \pg{35--43}}
\begin{fbx}[Extended G--K \pg{35}]
\[M=D(N-J-1)+\sum_{i=1}^{J}F_i+R_C-R_M,\qquad R_C=\sum_iF_i-\sum_{k=1}^{L}u_k\]
$R_C$: over-constraints (constraints repeated by several joints; add back). $R_M$: local mobility (internal DoF that do not move the output; subtract). $L$ loops; $u_k$ = platform DoF loop $k$ allows (union, shared DoF once).
\end{fbx}
\begin{recipe}[extended count]
\begin{enumerate}
\item Plain G--K; write the number.
\item $R_M$: a rod with S joints at both ends spins about its axis: $+1$ each (U ends: 0).
\item $R_C$: per loop list the platform DoF allowed; take the union; $R_C=\sum F_i-\sum u_k$.
\item $M=$ plain $+R_C-R_M$. Check against the motion you expect.
\end{enumerate}
\end{recipe}
\begin{fbx}[Plain vs extended]
\footnotesize
\begin{tabularx}{\linewidth}{@{}Xccccrccr@{}}\toprule
Mechanism & $D$ & $N$ & $J$ & $\sum F$ & plain & $R_C$ & $R_M$ & $M$\\\midrule
3-PRR \pg{27} & 3 & 8 & 9 & 9 & 3 & 0 & 0 & 3\\
Simplest delta \pg{30} & 6 & 8 & 9 & 15 & 3 & 0 & 0 & 3\\
Common delta \pg{33} & 6 & 11 & 15 & 39 & 9 & 0 & 6 & 3\\
4RRR \pg{38} & 6 & 10 & 12 & 12 & $-6$ & 7 & 0 & 1\\
3-RPS \added & 6 & 8 & 9 & 15 & 3 & 0 & 0 & 3\\
3-SPS \added & 6 & 8 & 9 & 21 & 9 & 0 & 3 & 6\\
Four-bar, slider-crank \added & 3 & 4 & 4 & 4 & 1 & 0 & 0 & 1\\
Planar five-bar \added & 3 & 5 & 5 & 5 & 2 & 0 & 0 & 2\\
Stewart 6-UPS \added & 6 & 14 & 18 & 36 & 6 & 0 & 0 & 6\\\bottomrule
\end{tabularx}
\end{fbx}
\begin{worked}[counting bodies and joints \chk]
3-PRR: ground, platform, 3 cylinders, 3 pistons ($N=8$); 3R + 3P + 3R ($J=9$).
Simplest delta: base, 3 arms, 3 rods, platform ($N=8$); 3R + 3U + 3U ($J=9$), $\sum F=3(1+2+2)$.
Common delta: base, 3 arms, 6 parallelogram bars, platform ($N=11$); 3R + 12S ($J=15$), $\sum F=3(1+4\cdot3)=39$. Plain $6(11-15-1)+39=9$, but the platform only translates: 6 bars spin about themselves, $R_M=6$, $R_C=0$: $\mathbf{M=3}$.
Stewart: ground, platform, $6\times2$ leg bodies ($N=14$); 6U + 6P + 6S ($J=18$).
4RRR: legs $4\times2$ links, 4 legs $\times$ 3R. Loop 1 allows $\{T_y,T_z,R_x\}$, loop 2 $\{T_x,T_z,R_y\}$ \pg{39--41}: union has 5, so $R_C=12-5=7$, $M=-6+7-0=1$ (platform slides along $z$ only).
\end{worked}
\begin{slideerr}[slides 29, 35--36]
Handwritten $D(N-J+1)$ on slide 29: wrong; the printed equation and all examples use $-1$ (the $+1$ would give 15 for the simplest delta). For the delta the slides just state $R_C=0$; the $u_k$ formula does not give that, so argue ``no repeated constraint''. \chk
\end{slideerr}
\begin{trap}
Treating G--K as exact: it assumes independent constraints. Equal parallel links (3 hanging links, $N=5$, $J=6$) give $M=0$ but the body really moves.
\end{trap}

\hd{Loop closure and four-bar \pg{45--58, 71}}
\premise[loop closure]{geometry and active joints $\boldsymbol\theta$ ($n$ of them)}{$m$ equations $\eta_i(\mathbf q)=0$ for the passive joints $\boldsymbol\varphi$}
\begin{recipe}[loop closure \pg{49}]
\begin{enumerate}
\item Cut into serial chains; pick the cut with the fewest passive unknowns.
\item For each chain write the pose of the cut point.
\item Equate the two descriptions at the cut.
\item Count: $n$ active, $m$ passive, $m$ equations. Four-bar: cut at K $\Rightarrow$ $n=1$, $m=3$; cut link 2 (U, V) $\Rightarrow$ $m=1$. Stewart: 6 active, 30 passive, not hand-solvable.
\end{enumerate}
\end{recipe}
\slidefig[0.92\linewidth]{w04-lec}{71}{0 119 538 151}{four-bar: {S} left pivot, {D} at $(L_0,0)$; K on link 2, $a,b$ its distances to the ends, $L_2=a+b$, output angle $\phi$}
\begin{fbx}[Four-bar, cut link 2 at K \pg{54}]
\[\begin{aligned}
L_1\cos\theta_1+a\cos(\theta_1+\varphi_2)&=L_0+L_3\cos\varphi_1+b\cos(\varphi_1+\varphi_3)\\
L_1\sin\theta_1+a\sin(\theta_1+\varphi_2)&=L_3\sin\varphi_1+b\sin(\varphi_1+\varphi_3)\\
\theta_1+\varphi_2&=\varphi_1+\varphi_3+\pi
\end{aligned}\]
$\theta_1$ active (crank angle at {S}); $\varphi_1$ at {D} to link 3; $\varphi_2,\varphi_3$ the coupler joints; all absolute, CCW, rad. Third equation: B$\to$K points back along the coupler (+$\pi$). \chk
\end{fbx}
\begin{fbx}[Freudenstein equation, cut both ends \pg{57}]
\[\eta_1=(L_1\cos\theta_1-L_0-L_3\cos\varphi_1)^2+(L_1\sin\theta_1-L_3\sin\varphi_1)^2-L_2^2=0\]
$n=m=1$. Given $\theta_1$: two roots $\varphi_1$ = two assembly modes. With $u=L_1\cos\theta_1-L_0$, $v=L_1\sin\theta_1$, $d^2=u^2+v^2$:
\[k=\frac{d^2+L_3^2-L_2^2}{2L_3},\quad\varphi_1=\operatorname{atan2}(v,u)\pm\arccos\frac kd\]
\[\varphi_2=\delta-\theta_1,\qquad\varphi_3=\theta_1+\varphi_2-\varphi_1-\pi\]
$\lvert k/d\rvert>1$: no assembly. ($\delta=\theta_1+\varphi_2$ = coupler direction.)
\end{fbx}

\hd{Four-bar IK (easy): K and $\phi$ given \pg{68--73}}
\premise[four-bar IK, W13 Q5]{K$=(x,y)$ in {S}, output angle $\phi$; $L_0,L_1,L_3,a,b$}{$\theta_1,\varphi_1,\varphi_2,\varphi_3$}
Easy because each leg is a serial chain with known end pose: solve legs separately, one answer per leg.
\begin{deriv}[\pg{70--72}]
Left {S}$\to$K: $x=L_1\cos\theta_1+a\cos\phi$, $y=L_1\sin\theta_1+a\sin\phi$ (using $\theta_1+\varphi_2=\phi$).
Right {D}$\to$K: $\varphi_1+\varphi_3=\phi+\pi$, $\cos(\varphi_1+\varphi_3)=-\cos\phi$: $L_3\cos\varphi_1=x-L_0+b\cos\phi$, $L_3\sin\varphi_1=y+b\sin\phi$.
\end{deriv}
\begin{recipe}[four-bar IK]
\begin{enumerate}
\item $\theta_1=\operatorname{atan2}(y-a\sin\phi,\ x-a\cos\phi)$.
\item $\varphi_2=\phi-\theta_1$ (add $2\pi$ if positive wanted).
\item $\varphi_1=\operatorname{atan2}(y+b\sin\phi,\ x+b\cos\phi-L_0)$.
\item $\varphi_3=\phi+\pi-\varphi_1$ (mod $2\pi$).
\item Check: $\lvert\mathbf A\rvert=L_1$, $\lvert\mathbf B-\mathbf D\rvert=L_3$, and $\theta_1+\varphi_2=\varphi_1+\varphi_3+\pi$.
\end{enumerate}
\end{recipe}
\begin{worked}[IK numbers]
$L_0=1.0$, $L_1=0.5$, $L_3=0.6$, $a=0.2$, $b=0.25$ m; K$=(0.3710,0.4678)$, $\phi=-0.600^\circ$.
$\theta_1=\operatorname{atan2}(0.4699,0.1710)=70.0^\circ$; $\varphi_2=289.4^\circ$ ($\equiv-70.6^\circ$); $\varphi_1=\operatorname{atan2}(0.4652,-0.3790)=129.17^\circ$; $\varphi_3=-0.6+180-129.175=50.2^\circ$. Check: $359.4^\circ=129.2+50.2+180$. \chk
\end{worked}
\begin{slideerr}[W13 p.\,11 (Q5)]
Printed $\varphi_3=\varphi_1-\phi+\pi$ puts K at $(0.669,0.711)$, not $(0.371,0.468)$. W4 p.\,72 ($\phi=\varphi_1+\varphi_3-\pi$) is right: use $\varphi_1+\varphi_3=\phi+\pi$. \chk
\end{slideerr}
\begin{trap}
Radians vs degrees; atan (not atan2) when $x+b\cos\phi-L_0<0$ (second quadrant); the printed $\varphi_3$.
\end{trap}
\textbf{W13 sample Q5} (four-bar IK), \textbf{Q6} (planar mobility), \textbf{Q7} (delta mobility): solved on p.\,\pageref{p:SAMPLE}.

\hd{Four-bar FK (hard) \pg{59--67}}
\premise[parallel FK]{$\theta_1$ and geometry}{$\varphi_i$, then the output pose; up to 2 assembly modes}
\textbf{Steps} \pg{60}: solve the non-linear loop equations for $\varphi_i$, then compute the pose. No general closed form for other parallel robots; numerics \pg{67}.
\begin{twoway}{A. Slide route (algebra) \pg{61--66}}{B. Geometry (circle intersection) \added}
Eliminate angles by squaring and adding: two equations $A_i\cos\delta+B_i\sin\delta+C_i=0$; Cramer or half-angle; then $\varphi_2=\delta-\theta_1$, $\varphi_1$ from Freudenstein.
\tcblower
B on circle $L_2$ about A \emph{and} circle $L_3$ about D; two intersections = two modes. Fastest by hand: use this unless told otherwise.
\end{twoway}
\begin{fbx}[Route B: geometric FK \added]
\[\mathbf A=L_1\begin{bmatrix}\cos\theta_1\\\sin\theta_1\end{bmatrix},\ \mathbf D=\begin{bmatrix}L_0\\0\end{bmatrix},\ d=\lVert\mathbf D-\mathbf A\rVert,\ \hat{\mathbf u}=\frac{\mathbf D-\mathbf A}{d},\ \hat{\mathbf n}=\begin{bmatrix}-u_y\\u_x\end{bmatrix}\]
\[x_a=\frac{L_2^2-L_3^2+d^2}{2d},\quad h=\sqrt{L_2^2-x_a^2},\quad\mathbf B=\mathbf A+x_a\hat{\mathbf u}\pm h\hat{\mathbf n}\]
\[\phi=\operatorname{atan2}(B_y-A_y,B_x-A_x),\ \ \varphi_2=\phi-\theta_1,\ \ \mathbf K=\mathbf A+a(\cos\phi,\sin\phi)\]
\[\varphi_1=\operatorname{atan2}(B_y,B_x-L_0),\quad\varphi_3=\phi+\pi-\varphi_1\ (\mathrm{mod}\ 2\pi)\]
$+$: B left of $\mathbf A\!\to\!\mathbf D$ (above); $-$: below. Needs $\lvert L_2-L_3\rvert\le d\le L_2+L_3$.
\end{fbx}
\begin{deriv}[slide route \pg{61--65}]
$x,y$ = K in {S}; $\delta=\theta_1+\varphi_2$; $\cos(\varphi_1+\varphi_3)=-\cos\delta$. Right chain: $L_3\cos\varphi_1=(x-L_0)+b\cos\delta$, $L_3\sin\varphi_1=y+b\sin\delta$; square, add:
$A_1=x-L_0,\ B_1=y,\ C_1=\frac1{2b}[(x-L_0)^2+y^2+b^2-L_3^2]$.
Left chain: $x-a\cos\delta=L_1\cos\theta_1$, $y-a\sin\delta=L_1\sin\theta_1$; square, add:
$A_2=x,\ B_2=y,\ C_2=\frac1{2a}[L_1^2-a^2-x^2-y^2]$.
Cramer, $\Delta=A_1B_2-A_2B_1$:
\[\cos\delta=\tfrac{B_1C_2-B_2C_1}{\Delta},\ \ \sin\delta=\tfrac{A_2C_1-A_1C_2}{\Delta},\ \ \delta=\operatorname{atan2}(\sin\delta,\cos\delta)\]
Half-angle $t=\tan\frac\delta2$: $(C_i-A_i)t^2+2B_it+(A_i+C_i)=0$; eliminating $t^2$:
\[t=\frac{A_2C_1-A_1C_2}{(B_1C_2-B_2C_1)+(A_1B_2-A_2B_1)},\qquad\delta=2\tan^{-1}t\ \ \chk\]
\end{deriv}
\begin{slideerr}[slide 65]
Printed $\delta=\tan^{-1}\frac{A_1C_2-A_2C_1}{(B_1C_2-B_2C_1)+(A_1B_2-A_2B_1)}$: sign flipped and factor 2 missing (mode $+$ test: slide $0.30^\circ$, true $-0.60^\circ$). Use Cramer/atan2. Slide's quartic holds for every K on the coupler curve and has no $\theta_1$. \chk
\end{slideerr}
\begin{worked}[four-bar FK, $\theta_1=70^\circ$, both modes]
Same lengths as the IK example; $L_2=0.45$. $\mathbf A=(0.1710,0.4698)$; $d^2=1.25-\cos70^\circ=0.9080$, $d=0.9529$; $\hat{\mathbf u}=(0.8700,-0.4931)$, $\hat{\mathbf n}=(0.4931,0.8700)$; $x_a=\frac{0.2025-0.36+0.9080}{1.9058}=0.3938$, $h=\sqrt{0.2025-0.1551}=0.2178$.
\begin{center}
\begin{tikzpicture}[x=6.4cm,y=5.4cm,line width=0.7pt]
\coordinate (S) at (0,0); \coordinate (D) at (1,0);
\coordinate (A) at (0.171,0.4698); \coordinate (Bp) at (0.621,0.4651); \coordinate (Bm) at (0.4062,0.0862);
\coordinate (Kp) at (0.371,0.4678); \coordinate (Km) at (0.2756,0.2993);
\draw[g3,line width=0.5pt] (S)--(D) node[midway,below,font=\scriptsize,text=black]{$L_0=1.0$};
\draw[line width=1.6pt] (S)--(A); \draw (A)--(Bp)--(D); \draw[dashed] (A)--(Bm)--(D);
\foreach \p in {S,D,A,Bp,Bm} \filldraw[fill=white] (\p) circle (0.014);
\fill (Kp) circle (0.01); \fill (Km) circle (0.01);
\node[font=\scriptsize,below left] at (S) {S}; \node[font=\scriptsize,below right] at (D) {D};
\node[font=\scriptsize,above left] at (A) {A}; \node[font=\scriptsize,above] at (Bp) {B$_{+}$};
\node[font=\scriptsize,below left] at (Bm) {B$_{-}$}; \node[font=\scriptsize,above right=0pt] at (Kp) {K$_+$};
\node[font=\scriptsize,right] at (Km) {K$_-$};
\draw (0.08,0) arc (0:70:0.08); \node[font=\scriptsize] at (0.135,0.045) {$\theta_1$};
\draw (1.12,0) arc (0:129.175:0.12); \node[font=\scriptsize] at (1.11,0.11) {$\varphi_1^{+}$};
\end{tikzpicture}
\end{center}
\drawncap{solid: mode $+$; dashed: mode $-$.}
\end{worked}
\begin{worked}[four-bar FK results \chk]
\begin{tabularx}{\linewidth}{@{}Xcc@{}}\toprule
 & mode $+$ & mode $-$\\\midrule
$\mathbf B$ (m) & $(0.6210,0.4651)$ & $(0.4062,0.0862)$\\
$\phi$ & $-0.600^\circ$ & $-58.486^\circ$\\
$\varphi_2$ & $289.400^\circ$ & $231.514^\circ$\\
$\varphi_1$ & $129.175^\circ$ & $171.739^\circ$\\
$\varphi_3$ & $50.225^\circ$ & $309.775^\circ$\\
$\mathbf K$ (m) & $(0.3710,0.4678)$ & $(0.2756,0.2993)$\\\bottomrule
\end{tabularx}
Freudenstein route: $u=-0.8290$, $v=0.4698$, $k=0.8879$, $\arccos(k/d)=21.282^\circ$, $\operatorname{atan2}(v,u)=150.457^\circ$: $\varphi_1=129.175^\circ$ or $171.739^\circ$. Closure (mode $+$): $0.1710+0.2\cos359.4^\circ=0.3710=1-0.3790+0.25\cos179.4^\circ$ \checkmark.
\end{worked}
\begin{recipe}[crank range and full turn \added]
\begin{enumerate}
\item Closes iff $\lvert L_2-L_3\rvert\le d\le L_2+L_3$, $d^2=L_0^2+L_1^2-2L_0L_1\cos\theta_1$.
\item $\dfrac{L_0^2+L_1^2-(L_2+L_3)^2}{2L_0L_1}\le\cos\theta_1\le\dfrac{L_0^2+L_1^2-(L_2-L_3)^2}{2L_0L_1}$; if it contains $[-1,1]$ the crank turns fully, else it rocks (modes merge at the limits).
\item Grashof: shortest + longest $\le$ other two $\Rightarrow$ link next to the shortest turns fully.
\end{enumerate}
\end{recipe}
Example above: $\cos\theta_1\ge0.1475$, $\lvert\theta_1\rvert\le81.52^\circ$; $0.45+1.0>0.5+0.6$: non-Grashof, crank rocks. \chk

\hd{Delta robot IK (physical branch) \added}
\premise[delta IK]{platform centre $\mathbf p$ in {B} ($z$ up, $p_z<0$); $R_b,R_p$, upper arm $L_a$, rod $L_r$}{$\theta_1,\theta_2,\theta_3$}
Platform stays level: solve leg $i$ alone, rod joint $\mathbf p+R_p\hat{\mathbf u}_i$.
\begin{deriv}[leg $i$, $\gamma_i=0^\circ,120^\circ,240^\circ$]
$\hat{\mathbf u}_i=(\cos\gamma_i,\sin\gamma_i,0)$; $\theta_i$ measured \emph{down} from horizontal; elbow $\mathbf E_i=R_b\hat{\mathbf u}_i+L_a(\cos\theta_i\hat{\mathbf u}_i-\sin\theta_i\hat{\mathbf z})$. With $\mathbf d_i=\mathbf p+(R_p-R_b)\hat{\mathbf u}_i$, rod length $L_r$:
\[L_a^2+\lVert\mathbf d_i\rVert^2-2L_a(d_u\cos\theta_i-d_z\sin\theta_i)=L_r^2,\ \ d_u=\mathbf d_i\!\cdot\!\hat{\mathbf u}_i,\ d_z=\mathbf d_i\!\cdot\!\hat{\mathbf z}\]
$A=d_u$, $B=-d_z$, $K=\dfrac{\lVert\mathbf d_i\rVert^2+L_a^2-L_r^2}{2L_a}$, $\rho=\sqrt{A^2+B^2}$:
\[A\cos\theta_i+B\sin\theta_i=K\ \Rightarrow\ \theta_i=\operatorname{atan2}(B,A)-\arccos(K/\rho),\ \ \lvert K\rvert\le\rho\]
\end{deriv}
\textbf{Physical root}: the minus root. The $+$ root folds the arm back above the base (elbow inside base circle, $z>0$). Pick elbow outside the base circle and below it: $\lVert\mathbf E_{xy}\rVert>R_b$, $E_z<0$.
\begin{worked}[delta numbers]
$R_b=0.20$, $R_p=0.05$, $L_a=0.20$, $L_r=0.45$ m; $\mathbf p=(0.03,-0.02,-0.38)$. Leg 1: $\mathbf d_1=(-0.12,-0.02,-0.38)$, $\lVert\mathbf d_1\rVert^2=0.1592$, $K=-0.00825$, $A=-0.12$, $B=0.38$, $\rho=0.3985$; $\operatorname{atan2}(B,A)=107.526^\circ$, $\arccos(-0.0207)=91.186^\circ$: $\theta_1=16.339^\circ$ (other root $198.7^\circ$ unphysical). $\theta_2=30.871^\circ$, $\theta_3=22.996^\circ$; rods all $0.45$ m. Home $\mathbf p=(0,0,-0.4)$: $\theta_i=27.277^\circ$. \chk
\end{worked}

\hd{Five-bar IK (2 DoF) \added}
\premise[five-bar IK]{tip $\mathbf P$; pivots $\mathbf O_1,\mathbf O_2$; proximal $L_p$, distal $L_d$}{$\theta_1,\theta_2$}
Each side is a 2R problem: $\mathbf v_j=\mathbf P-\mathbf O_j$, $r_j=\lVert\mathbf v_j\rVert$, cosine rule in $\mathbf O_j\mathbf E_j\mathbf P$:
\[\cos\beta_j=\frac{r_j^2+L_p^2-L_d^2}{2r_jL_p},\qquad\theta_j=\operatorname{atan2}(v_{j,y},v_{j,x})\pm\beta_j\]
Reachable iff $\lvert L_p-L_d\rvert\le r_j\le L_p+L_d$; 4 configurations, ``elbows outward'': $+\beta_1,-\beta_2$.
\begin{worked}[five-bar numbers]
$\mathbf O_{1,2}=(\mp0.1,0)$, $L_p=0.3$, $L_d=0.4$; $\mathbf P=(0.05,0.45)$. Left: $r_1=0.4743$, $\cos\beta_1=0.5446$, $\beta_1=57.00^\circ$, $\operatorname{atan2}=71.565^\circ$: $\theta_1=128.567^\circ$. Right: $r_2=0.4528$, $\beta_2=60.20^\circ$, $\operatorname{atan2}=96.340^\circ$: $\theta_2=36.138^\circ$. Check $\lVert\mathbf P-\mathbf E_j\rVert=0.4$. FK: intersect the two $L_d$ circles about $\mathbf E_1,\mathbf E_2$ (2 points). \chk
\end{worked}
\textbf{Singularities} \added: none on slides (W6: $\det\mathbf J=0$). Four-bar modes merge at $d=L_2\pm L_3$ (dead centre); a leg is singular stretched/folded ($r_j=L_p\pm L_d$).

\hd{Slider-crank and Stewart leg IK \added}
\begin{fbx}[Slider-crank (crank $L_1$, rod $L_2>L_1$, slider on the pivot line)]
\[x=L_1\cos\theta+\sqrt{L_2^2-L_1^2\sin^2\theta},\qquad\beta=-\arcsin\frac{L_1\sin\theta}{L_2}\]
Closure: tip height $L_1\sin\theta=-L_2\sin\beta$. $L_1=0.1$, $L_2=0.3$: $x(0^\circ)=0.4000$, $x(60^\circ)=0.3372$, $x(90^\circ)=0.2828$, $x(180^\circ)=0.2$; stroke $2L_1$. \chk
\end{fbx}
\begin{fbx}[Stewart leg $i$ \pg{73}; formulas \added]
\[{}^B\mathbf l_i=\mathbf t+\Rm\mathbf p_i-\mathbf b_i,\quad d_i=\lVert{}^B\mathbf l_i\rVert,\quad{}^B\hat{\mathbf n}_i=\frac{{}^B\mathbf l_i}{d_i},\quad{}^P\hat{\mathbf n}_i=\Rm^{\mathsf T}\,{}^B\hat{\mathbf n}_i\]
\[\psi_i=\operatorname{atan2}(n_y,n_x),\quad\varepsilon_i=\arcsin n_z\]
$\mathbf b_i$ in {B}, $\mathbf p_i$ in {P}, $\mathbf t,\Rm$ pose of {P} in {B}; $d_i$ active length. U-axis convention not fixed by the slide: state yours.
Base circle 1.0 m ($-10,10,110,130,230,250^\circ$), platform circle 0.5 m ($-50,50,70,170,190,290^\circ$), home $\mathbf t=(0,0,1)$: $\mathbf l_1=(-0.6634,-0.2094,1.0)$, $d_1=1.2182$ m (all six equal). \chk
\end{fbx}

\begin{examq}[assembly modes, crank test, mobility \mk{marks assumed}]
Four-bar $L_0=0.8$, crank $L_1=0.3$ at {S}$=(0,0)$, coupler $L_2=0.7$, rocker $L_3=0.5$ at $\mathbf D=(0.8,0)$ m. (a) At $\theta_1=60^\circ$ find both B and $\phi$. (b) Does the crank turn fully? (c) $M$?
\solution
(a) $\mathbf A=(0.1500,0.2598)$; $d^2=0.64+0.09-0.24=0.49$, $d=0.7$; $x_a=\frac{0.49-0.25+0.49}{1.4}=0.5214$, $h=\sqrt{0.49-0.2719}=0.4670$, $\hat{\mathbf u}=(0.9286,-0.3712)$, $\hat{\mathbf n}=(0.3712,0.9286)$. B$_+=\boxed{(0.8075,0.4999)}$, $\phi=20.06^\circ$; B$_-=\boxed{(0.4608,-0.3674)}$, $\phi=-63.64^\circ$.\\
(b) $0.3+0.8=1.1\le0.5+0.7=1.2$, shortest is the crank: crank-rocker, turns fully ($d$ runs $0.5$ to $1.1\in[0.2,1.2]$).\\
(c) $3(4-4-1)+4=\boxed1$.
\end{examq}
\begin{examq}[spot the errors: simplest delta \mk{marks assumed}]
``$D=6$, $N=8$, $J=9$, $\sum F=15$, $M=6(8-9+1)+15=15$. The rod spins, $R_M=6$, so $M=9$.'' Find both errors.
\solution
(1) Sign: $D(N-J-\mathbf1)$: $M=-12+15=\boxed3$. (2) Rods have a U joint at each end; U transmits spin about the rod axis, so $R_M=0$; the student invented $R_M$ to patch error 1. Contrast common delta (ball joints): plain 9, $R_M=6$, true 3.
\end{examq}
\begin{examq}[planar 3-RPR IK \mk{marks assumed}]
$\mathbf b_1=(0,0)$, $\mathbf b_2=(1.0,0)$, $\mathbf b_3=(0.5,0.9)$ m; platform points $\mathbf a_i=0.2(\cos\gamma_i,\sin\gamma_i)$, $\gamma_i=0^\circ,120^\circ,240^\circ$. Leg lengths at $(x,y,\varphi)=(0.5,0.35,10^\circ)$? Why is FK harder?
\solution
$\mathbf p_i=\mathbf t+\Rm(\varphi)\mathbf a_i$, $\ell_i=\lVert\mathbf p_i-\mathbf b_i\rVert$: $\mathbf p_1=(0.6970,0.3847)$, $\mathbf p_2=(0.3714,0.5032)$, $\mathbf p_3=(0.4316,0.1621)$.
\[\boxed{\ell_1=0.7961,\ \ell_2=0.8052,\ \ell_3=0.7411\ \text{m}}\]
Wrong: not rotating $\mathbf a_i$ ($\ell_1=0.7826$). FK: three coupled non-linear loop equations, several solutions; IK: none. \chk
\end{examq}
\end{multicols}
